\section*{Bab \ref{ch:b3:astrophysics} --- Astrofisika}

\begin{solution}{exo:b3:astrophysics:1}
(a) $P_{\text{c}} \sim GM^2/R^4 \approx \qty{1.1e15}{Pa}$ --- yakni sepuluh
miliar atmosfer. (b) $\bar\rho \approx \qty{1400}{kg/m^3}$:
$T \sim P m_{\text{p}}/\rho k_{\text{B}} \approx \qty{e8}{K}$ ---
kasar, satu orde di atas nilai sejatinya \qty{1.6e7}{K}, tetapi skalanya
tepat. (c) Persis pada ambang penerowongannya. (d) Sebuah protobintang
mengerut dan memanas \emph{sampai} terasnya menyala, lalu berhenti
mengerut: bola gas mana pun yang cukup masif pasti tiba pada penyalaannya
--- bintangnya menyetel dirinya pada termostat intinya.
\end{solution}

\begin{solution}{exo:b3:astrophysics:2}
(a) $L = 4\pi d^2\times1360 \approx \qty{3.8e26}{W}$. (b) $T =
(L/4\pi R^2\sigma)^{1/4} \approx \qty{5800}{K}$. (c) Wien:
$\lambda_{\text{puncak}} \approx \qty{500}{nm}$ --- matanya berevolusi
tepat pada puncak Matahari. (d) Fluks kuadrat terbaliknya (Tahun 1);
Stefan--Boltzmann dan Wien: yakni benda hitam \cref{ch:b3:photons-phonons},
yang diturunkan dari statistika foton.
\end{solution}

\begin{solution}{exo:b3:astrophysics:3}
(a) Sirius B: jauh di kiri, jauh di bawah; Betelgeuse: jauh di kanan, jauh di atas.
(b) $R/R_\odot = \sqrt{L/L_\odot}\,(T_\odot/T)^2$: Sirius B
$\approx 0.009\,R_\odot \approx$ satu jejari Bumi; Betelgeuse
$\approx 860\,R_\odot \approx \qty{4}{au}$. (c) Merkurius, Venus,
Bumi dan Mars akan mengorbit \emph{di dalamnya}. (d) Seukuran Bumi namun
bermassa matahari: hanya \omterm{thm:b3:quantum-statistics:fermi}{tekanan kemerosotan} elektron
(\cref{ch:b3:quantum-statistics}) --- Sirius B adalah
\omterm{ex:b3:quantum-statistics:dwarf}{katai putih} yang klasik.
\end{solution}

\begin{solution}{exo:b3:astrophysics:4}
(a) $10\,\text{Gyr}\times10^{-2.5} \approx \qty{30}{Myr}$. (b)
$0.5^{-2.5} \approx 5.7$: yakni kira-kira 57 miliar tahun --- tiap katai
merah yang pernah lahir masih membakar. (c) Kecerlangannya tumbuh sebagai
$M^{3.5}$ tetapi tangkinya hanya sebagai $M$: si kaya membakar hartanya
lebih cepat daripada menabungnya. (d) Bintang masif mati dalam beberapa
juta tahun sejak kelahirannya --- sebelum menghanyut ke mana pun: supernova
meledak di dalam lengan spiral pembentuk bintang yang membuatnya.
\end{solution}

\begin{solution}{exo:b3:astrophysics:5}
(a) $E = -K$: memancarkan membuat $E$ makin negatif, sehingga $K$ ---
yakni suhunya --- \emph{menaik}: yakni tanda tangan termodinamika gravitasi. (b) $GM^2/R \approx \qty{3.8e41}{J}$. (c)
$E/L_\odot \approx \qty{e15}{s} \approx \qty{30}{Myr}$. (d)
Matahari Kelvin tak mungkin lebih tua daripada puluhan megatahun, sedangkan
cacing Darwin dan batuan Lyell menuntut miliaran: yakni
penyelesaiannya adalah bahan bakar yang tak dikenal siapa pun --- \cref{ch:b3:nuclear-physics},
tujuh puluh tahun terlalu awal.
\end{solution}

\begin{solution}{exo:b3:astrophysics:6}
(a) $\bar\rho \approx \qty{1.8e9}{kg/m^3}$: yakni sendok teh
sembilan ton. (b) $g = GM/R^2 \approx \qty{3.3e6}{m/s^2}$ --- yakni tiga
ratus ribu gravitasi Bumi. (c) $n \propto M/R^3$:
kemerosotannya memasok $P \propto M^{5/3}/R^5$, gravitasinya menuntut
$P \propto M^2/R^4$; menyamakannya, $R \propto M^{-1/3}$. (d) Berat
yang lebih besar memerlukan tekanan Fermi yang lebih besar, yang hanya disediakan
pemampatannya --- sehingga massa justru \emph{menyusutkan} bintangnya; pusarannya punya
dinding: begitu elektronnya menjadi relativistik hukum tekanannya melunak menjadi
$n^{4/3}$, keseimbangannya gagal pada jejari berapa pun, dan
$M_{\text{Ch}} = 1.4\,M_\odot$ adalah tepi tebingnya.
\end{solution}

\begin{solution}{exo:b3:astrophysics:7}
(a) $V = M/\rho_{\text{inti}} \approx \qty{1.2e13}{m^3}$:
$R \approx \qty{14}{km}$ --- yakni sebuah inti bercakrawala kota. (b)
$\omega \propto 1/R^2$: percepatan putarnya $(10^8/1.4\times10^4)^2
\approx 5\times10^{7}$, sehingga 25 hari menyusut menjadi
${\sim}\qty{40}{ms}$: yakni wilayah pulsar. (c) $g \approx
\qty{e12}{m/s^2}$: sebuah gunung di sini akan setinggi milimeter di sana.
(d) $B \approx 10^{-2}\times5\times10^{7} \approx
\qty{5e5}{T}$ --- dan sumbu magnetnya, yang tak sejajar dengan sumbu
putarnya, menyalurkan pancaran radio sepanjang kutubnya: berkasnya
menyapu seperti mercusuar, dan kilatannya kita sebut sebuah pulsar.
\end{solution}

\begin{solution}{exo:b3:astrophysics:8}
(a) $\tfrac12mc^2 = GMm/R$ memberikan $R = 2GM/c^2$. (b) Matahari:
\qty{3.0}{km}; Bumi: \qty{8.9}{mm} --- mampatkanlah salah satunya ke dalam
jejari itu dan cahaya tak dapat pergi. (c) $R_{\text{s}} \approx
\qty{1.2e10}{m}$: yakni seperlima orbit Merkurius --- dan bintang teramati
berputar kencang mengelilinginya. (d) $\bar\rho \propto
M/R_{\text{s}}^3 \propto 1/M^2$: lubang bermassa $10^8\,M_\odot$ kurang
rapat daripada air. Tak ada yang perlu ``diremukkan'' pada cakrawalanya ---
ia adalah titik tanpa jalan kembali, bukan permukaan batu.
\end{solution}

\begin{solution}{exo:b3:astrophysics:9}
(a) $H_0 = 7\times10^4/3.09\times10^{22} \approx
\qty{2.3e-18}{s^{-1}}$. (b) $1/H_0 \approx \qty{4.4e17}{s}
\approx 14$ miliar tahun. (c) $v = \qty{1.4e4}{km/s}$;
$z \approx v/c \approx 0.047$: panjang gelombangnya tiba terentang
5\,\%. (d) Lajunya tak pernah tetap --- diperlambat oleh
materi pada awalnya, lalu dipercepat oleh energi gelap sejak itu: bahwa
$1/H_0$ toh mendarat dalam satu gigatahun dari nilai sejatinya 13,8 adalah
kebetulan kosmik yang kecil.
\end{solution}

\begin{solution}{exo:b3:astrophysics:10}
(a) $\lambda_{\text{puncak}} = b/T \approx \qty{1.1}{mm}$:
yakni gelombang mikro, sesuai iklannya. (b) Terentang sebesar $3000/2.725 \approx
1100$ --- alam semesta seribu seratus kali lebih besar pada tiap
arahnya sejak itu. Penggabungan kembalinya menunggu pada \qty{3000}{K}, jauh
di bawah $\qty{13.6}{eV}/k_{\text{B}}$, karena keseimbangan Saha
(\cref{prop:b3:grand-canonical:saha}) dimiringkan oleh dua miliar
foton tiap \omterm{def:b3:particle-physics:zoo}{barionnya}: ekor pengionnya menjaga hidrogennya tetap terurai
lama setelah energetika naifnya berkata sebaliknya. (c) Kira-kira
$1.6\times10^{9}$ foton tiap \omterm{def:b3:particle-physics:zoo}{barionnya} --- alam semesta, menurut
cacahnya, nyaris seluruhnya cahaya. (d) CMB bukanlah sebuah tempat melainkan
seluruh langit: tiap garis pandangnya berakhir pada kabut hamburan
terakhirnya, sehingga tiap antena, yang diarahkan ke mana pun, menerimanya.
\end{solution}

\begin{solution}{exo:b3:astrophysics:11}
(a) $n/p = \eu^{-1.29/0.8} \approx 0.20$. (b) Dengan $n/p = 1/7$:
14 proton tiap 2 neutronnya; kedua neutronnya menyambar 2 proton ke dalam
satu $^4$He (bermassa 4) lalu menyisakan 12 proton sendirian: $4/16 = 25\,\%$ menurut
massanya. (c) Tak ada bintang yang punya waktu untuk menjadikan seperempat segalanya
helium; menemukan lantai itu pada gas tertua yang paling miskin logam
adalah alam semesta dini yang panas yang tertangkap basah --- yakni sebuah faktor
Boltzmann yang tertulis melintasi langit. (d) Tak ada inti yang mantap
pada $A = 5$ atau $8$, dan kerapatan serta suhunya sedang
turun dalam hitungan menit: jembatan menuju karbonnya (yakni tiga helium
sekaligus) hanya terbuka dalam teras \omterm{rem:b3:astrophysics:evolution}{raksasa merah} yang rapat dan sabar.
\end{solution}

\begin{solution}{exo:b3:astrophysics:12}
(a) $mv^2/r = GM(r)m/r^2$. (b) $M = v^2r/G \approx
\qty{1.1e42}{kg} \approx 5\times10^{11}\,M_\odot$. (c) Kira-kira
sepersepuluh: sembilan puluh persen galaksinya tak bersinar dan tak pula
menyerap. (d) $v$ yang datar berarti $M(r) \propto r$ --- massa yang tak
terlihat itu terus menumpuk jauh di balik cakram kasat matanya, yakni sebuah
lingkung gelap yang meluas (dengan kerapatan $\sim1/r^2$) yang di dalamnya
galaksi yang bercahaya hanyalah teras yang menyala.
\end{solution}

\begin{solution}{pb:b3:astrophysics:1}
\textbf{1.} $t \sim 10\,\text{Gyr}\times20^{-2.5} \approx
\qty{6}{Myr}$: Matahari akan hidup dua ribu kali umur itu.
\textbf{2.} Inti yang lebih berat membawa muatan lebih besar: penghalang
Coulomb yang lebih tinggi memerlukan teras yang lebih panas agar peluang
penerowongannya (\cref{exo:b3:nuclear-physics:8}) tetap mengimbangi.
\textbf{3.} Besi duduk di puncak kurva $B/A$: memfusikannya justru
\emph{berongkos} energi --- tangga bahan bakar tanurnya berakhir pada abu.
\textbf{4.} Yakni \omterm{prop:b3:astrophysics:compact}{massa Chandrasekhar}: yaitu paling banyak yang dapat
diangkut \omterm{thm:b3:quantum-statistics:fermi}{tekanan kemerosotan} elektron; terasnya adalah \omterm{ex:b3:quantum-statistics:dwarf}{katai putih} yang tumbuh
di dalam bintang yang hidup, dan pada $1.4\,M_\odot$ topangannya gagal.
\textbf{5.} Hidrogen: megatahun; helium: ratusan milenium;
silikon: satu hari --- tiap tahapnya membayar lebih sedikit tiap kilogramnya sementara
bintangnya, yang kian panas, membelanjakannya (sebagian besar ke neutrino) kian cepat:
jamnya berlari gila secara logaritmik.
\textbf{6.} $t \sim 1/\sqrt{G\rho} =
1/\sqrt{6.7\times10^{-11}\times10^{12}} \approx \qty{0.1}{s}$:
terasnya jatuh dari bawah bintangnya dalam sepersepuluh detik.
\textbf{7.} $E \sim GM^2/R \approx \qty{4e46}{J}$.
\textbf{8.} Pertunjukan fotonnya adalah $10^9\times3.8\times10^{26}
\times2.6\times10^{6} \approx \qty{e42}{J}$ --- yakni satu bagian dalam
sepuluh ribu. Sisanya pergi sebagai neutrino: materi inti yang runtuh
bersifat kedap terhadap segala hal lainnya, tetapi
(\cref{ch:b3:particle-physics}) nyaris tembus terhadap partikel gaya
lemahnya sendiri, yang menguras energinya dalam hitungan detik.
\textbf{9.} $E/4\pi d^2 = 4\times10^{46}/2.8\times10^{43}
\approx \qty{1.4e3}{J/m^2}$.
\textbf{10.} Pada $\qty{15}{MeV} = \qty{2.4e-12}{J}$ masing-masing:
${\sim}6\times10^{14}$ neutrino tiap meter persegi --- menembus
tiap meter persegi Bumi.
\textbf{11.} $6\times10^{14}\times10^{-46}\times1.5\times
10^{32} \approx 9$: pendeteksinya mencatat sebelas --- yakni
sasaran tepat berorde besar bagi bintang yang mati 160.000 tahun
lalu.
\textbf{12.} Neutrinonya meninggalkan terasnya pada keruntuhannya;
cahayanya baru lahir ketika kejutnya menembus selubung
bintangnya beberapa jam kemudian. Bahwa keduanya tiba pada malam yang sama setelah
160 milenium dalam penerbangan membatasi neutrinonya pada laju cahaya
dalam beberapa bagian dari $10^{9}$ --- dan massanya menjadi nyaris tak ada.
\textbf{13.} Anggaran energinya ($\sim10^{46}\,\unit{J}$, dengan 99\,\%
dalam neutrino) dan lama letusannya ${\sim}10$ detik ---
keduanya cocok dengan teras runtuh teorinya, sehingga menaikkan
mekanisme supernovanya dari papan tulis menjadi pengamatan.
\textbf{14.} $\bar\rho = 2.8\times10^{30}/(\tfrac43\pi
(1.2\times10^4)^3) \approx \qty{4e17}{kg/m^3}$: yakni kerapatan
inti --- yaitu sendok teh \cref{def:b3:nuclear-physics:nuclide},
yang kini seukuran kota.
\textbf{15.} Percepatan putarnya $(10^8/1.2\times10^4)^2 \approx
7\times10^{7}$: 30 hari $\to$ ${\sim}\qty{40}{ms}$.
\textbf{16.} $B \approx 10^{-2}\times7\times10^{7} \approx
\qty{7e5}{T}$ --- yakni seratus juta tesla-meter fluks,
yang dikekalkan ke dalam sebuah perangko.
\textbf{17.} Sumbu magnetnya, yang miring dari sumbu putarnya,
memancarkan berkas radio sepanjang kutubnya; berkasnya menyapu langit seperti
mercusuar. Sisa Kepitingnya --- yakni bintang tamu tahun 1054 --- masih
berputar tiga puluh kali tiap detiknya, dengan perlambatan putar semilenium yang
baru saja bermula.
\textbf{18.} $R_{\text{s}} = 3\times\qty{3}{km} \approx
\qty{9}{km}$ --- yakni seorde dengan bintangnya sendiri: tanpa
tekanan tersisa untuk berdebat, keruntuhannya berlanjut menembus cakrawalanya
dan sisanya menjadi sebuah \omterm{prop:b3:astrophysics:compact}{lubang hitam}.
\textbf{19.} Larik jam semacam itu adalah pendeteksi gelombang gravitasi
berbukaan galaksi (sekaligus laboratorium terbaik kenisbian:
orbit pulsar ganda meluruh persis seperti yang diramalkan radiasi
gravitasi).
\textbf{20.} Oksigennya difusikan dalam kulit sebuah bintang masif
lalu dilemparkan keluar oleh sebuah supernova; besinya ditempa dalam
akhir yang meledak itu sendiri; keduanya menghanyut dalam awan antarbintang,
bergabung dengan nebula suryanya, Bumi, lautnya, rantai makanannya
--- dan sang pembaca. Hidrogen dentuman besarnya, sebaliknya: itu sudah
menjadi milik Anda sejak dahulu.
\textbf{21.} Massa peledakan yang sama $\to$ kecerlangan hakiki
yang sama $\to$ kecerlangan tampaknya membaca jaraknya: yakni sebuah lilin
berdaya yang diketahui yang kasat mata sepanjang miliaran tahun cahaya ---
yaitu penggaris yang menggambar ujung jauh diagram Hubble.
\textbf{22.} Pemuaiannya \emph{memecut}: sesuatu yang disebut
``energi gelap'' mendorong alam semesta berpisah --- yakni Nobel 2011,
dan tagihan terbuka terbesar fisika.
\textbf{23.} $\qty{50}{kpc} \approx 163000$ tahun cahaya: yakni
cahayanya berangkat ketika \emph{Homo sapiens} baru meninggalkan Afrika.
\textbf{24.} $6\times10^{14}\times0.03 \approx 2\times10^{13}$
menembus tiap manusia yang hidup --- dan, menurut peluang
\cref{exo:b3:particle-physics:10}, tak satu pun dari sejuta orang
berinteraksi bahkan dengan satu neutrino: kilatan yang melebihi kecerlangan sebuah
galaksi melintasi umat manusia tanpa terasa.
\textbf{25.} Sebuah bintang yang seimbang sepuluh juta tahun di atas proton
yang menerowong; sebuah keruntuhan sepersepuluh detik yang membayarkan $10^{46}$
joule dalam neutrino; asas Pauli yang menahan jenazahnya pada
kerapatan inti; dan puing yang terserak --- oksigen, besi,
sang pembaca --- masih menunggangi pemuaian Hubble, dalam alam semesta yang
foto bayinya adalah sebuah kurva Planck. Fisika, lengkap untuk saat ini:
jilidnya ditutup.
\end{solution}
